%
\documentclass[12pt]{article}

\usepackage[usenames]{color}
\usepackage{verbatim}
\usepackage{fancyhdr}
\usepackage{graphicx}
\usepackage{amssymb,amsthm}
\usepackage{changepage}
%\usepackage[fleqn]{amsmath}
\usepackage{amsmath}
% \usepackage[final,notref,notcite,color]{showkeys}
\usepackage[left=1in,top=1in,right=1in,bottom=1in,letterpaper]{geometry}
\renewcommand{\baselinestretch}{1.2}

\pagestyle{fancy}
\lhead{}
\chead{}
\rhead{}
\lfoot{}
\cfoot{\thepage}
\rfoot{}



%% editing
\usepackage[shortlabels]{enumitem}
\usepackage[normalem]{ulem}

%% macros for commenting
\usepackage{mathtools}
\usepackage[normalem]{ulem} % to use \sout

%% math packages
% \usepackage{amsmath} % not needed for Springer
\usepackage{amssymb}
\usepackage{mathrsfs}

%\usepackage{subfigure}
%\usepackage{url}

%% a light-weight algorithm environment
\newtheorem{algo}{Algorithm}

%% highlighting and commenting
\newcommand{\outline}[1]{{\color{brown}#1}}
\newcommand{\rev}[1]{{\color{blue}#1}}
\newcommand{\remove}[1]{{\sout{#1}}}
\newcommand{\cut}[1]{{}}

%% macros for letters

\newcommand{\va}{{\mathbf{a}}}
\newcommand{\vb}{{\mathbf{b}}}
\newcommand{\vc}{{\mathbf{c}}}
\newcommand{\vd}{{\mathbf{d}}}
\newcommand{\ve}{{\mathbf{e}}}
\newcommand{\vf}{{\mathbf{f}}}
\newcommand{\vg}{{\mathbf{g}}}
\newcommand{\vh}{{\mathbf{h}}}
\newcommand{\vi}{{\mathbf{i}}}
\newcommand{\vj}{{\mathbf{j}}}
\newcommand{\vk}{{\mathbf{k}}}
\newcommand{\vl}{{\mathbf{l}}}
\newcommand{\vm}{{\mathbf{m}}}
\newcommand{\vn}{{\mathbf{n}}}
\newcommand{\vo}{{\mathbf{o}}}
\newcommand{\vp}{{\mathbf{p}}}
\newcommand{\vq}{{\mathbf{q}}}
\newcommand{\vr}{{\mathbf{r}}}
\newcommand{\vs}{{\mathbf{s}}}
\newcommand{\vt}{{\mathbf{t}}}
\newcommand{\vu}{{\mathbf{u}}}
\newcommand{\vv}{{\mathbf{v}}}
\newcommand{\vw}{{\mathbf{w}}}
\newcommand{\vx}{{\mathbf{x}}}
\newcommand{\vy}{{\mathbf{y}}}
\newcommand{\vz}{{\mathbf{z}}}

\newcommand{\ta}{{\tilde{a}}}
\newcommand{\tb}{{\tilde{b}}}
\newcommand{\tc}{{\tilde{c}}}
\newcommand{\td}{{\tilde{d}}}
\newcommand{\te}{{\tilde{e}}}
\newcommand{\tf}{{\tilde{f}}}
\newcommand{\tg}{{\tilde{g}}}

\newcommand{\ti}{{\tilde{i}}}
\newcommand{\tj}{{\tilde{j}}}
\newcommand{\tk}{{\tilde{k}}}
\newcommand{\tl}{{\tilde{l}}}
\newcommand{\tm}{{\tilde{m}}}
\newcommand{\tn}{{\tilde{n}}}

\newcommand{\tp}{{\tilde{p}}}
\newcommand{\tq}{{\tilde{q}}}

\newcommand{\ts}{{\tilde{s}}}

\newcommand{\tu}{{\tilde{u}}}
\newcommand{\tv}{{\tilde{v}}}
\newcommand{\tw}{{\tilde{w}}}
\newcommand{\tx}{{\tilde{x}}}
\newcommand{\ty}{{\tilde{y}}}
\newcommand{\tz}{{\tilde{z}}}

\newcommand{\vA}{{\mathbf{A}}}
\newcommand{\vB}{{\mathbf{B}}}
\newcommand{\vC}{{\mathbf{C}}}
\newcommand{\vD}{{\mathbf{D}}}
\newcommand{\vE}{{\mathbf{E}}}
\newcommand{\vF}{{\mathbf{F}}}
\newcommand{\vG}{{\mathbf{G}}}
\newcommand{\vH}{{\mathbf{H}}}
\newcommand{\vI}{{\mathbf{I}}}
\newcommand{\vJ}{{\mathbf{J}}}
\newcommand{\vK}{{\mathbf{K}}}
\newcommand{\vL}{{\mathbf{L}}}
\newcommand{\vM}{{\mathbf{M}}}
\newcommand{\vN}{{\mathbf{N}}}
\newcommand{\vO}{{\mathbf{O}}}
\newcommand{\vP}{{\mathbf{P}}}
\newcommand{\vQ}{{\mathbf{Q}}}
\newcommand{\vR}{{\mathbf{R}}}
\newcommand{\vS}{{\mathbf{S}}}
\newcommand{\vT}{{\mathbf{T}}}
\newcommand{\vU}{{\mathbf{U}}}
\newcommand{\vV}{{\mathbf{V}}}
\newcommand{\vW}{{\mathbf{W}}}
\newcommand{\vX}{{\mathbf{X}}}
\newcommand{\vY}{{\mathbf{Y}}}
\newcommand{\vZ}{{\mathbf{Z}}}

\newcommand{\cA}{{\mathcal{A}}}
\newcommand{\cB}{{\mathcal{B}}}
\newcommand{\cC}{{\mathcal{C}}}
\newcommand{\cD}{{\mathcal{D}}}
\newcommand{\cE}{{\mathcal{E}}}
\newcommand{\cF}{{\mathcal{F}}}
\newcommand{\cG}{{\mathcal{G}}}
\newcommand{\cH}{{\mathcal{H}}}
\newcommand{\cI}{{\mathcal{I}}}
\newcommand{\cJ}{{\mathcal{J}}}
\newcommand{\cK}{{\mathcal{K}}}
\newcommand{\cL}{{\mathcal{L}}}
\newcommand{\cM}{{\mathcal{M}}}
\newcommand{\cN}{{\mathcal{N}}}
\newcommand{\cO}{{\mathcal{O}}}
\newcommand{\cP}{{\mathcal{P}}}
\newcommand{\cQ}{{\mathcal{Q}}}
\newcommand{\cR}{{\mathcal{R}}}
\newcommand{\cS}{{\mathcal{S}}}
\newcommand{\cT}{{\mathcal{T}}}
\newcommand{\cU}{{\mathcal{U}}}
\newcommand{\cV}{{\mathcal{V}}}
\newcommand{\cW}{{\mathcal{W}}}
\newcommand{\cX}{{\mathcal{X}}}
\newcommand{\cY}{{\mathcal{Y}}}
\newcommand{\cZ}{{\mathcal{Z}}}

\newcommand{\ri}{{\mathrm{i}}}
\newcommand{\rr}{{\mathrm{r}}}

\newcommand{\EE}{{\mathbb{E}}}


%% macros for math notions and operators

\newcommand{\RR}{\mathbb{R}}
\newcommand{\CC}{\mathbb{C}}
\newcommand{\ZZ}{\mathbb{Z}}
\renewcommand{\SS}{{\mathbb{S}}}
\newcommand{\SSp}{\mathbb{S}_{+}}
\newcommand{\SSpp}{\mathbb{S}_{++}}
\newcommand{\sign}{\mathrm{sign}}
\newcommand{\vzero}{\mathbf{0}}
\newcommand{\vone}{{\mathbf{1}}}

\newcommand{\st}{{\text{s.t.}}} % subject to
\newcommand{\St}{{\mathrm{subject~to}}} % subject to
\newcommand{\op}{{\mathrm{op}}} % subscript for operator norm
\newcommand{\opt}{{\mathrm{opt}}} % subscript for optimal solution
%\newcommand{\supp}{{\mathrm{supp}}} % support
\newcommand{\Prob}{{\mathrm{Prob}}} % probability
\newcommand{\Diag}{{\mathrm{Diag}}} % vector -> diagonal matrix
%\newcommand{\diag}{{\mathrm{diag}}} % matrix diagonal -> vector
\newcommand{\dom}{{\mathrm{dom}}} % domain
\newcommand{\range}{{\mathrm{range}}} % domain
%\newcommand{\grad}{{\nabla}}    % gradient
\newcommand{\tr}{{\mathrm{tr}}} % trace
\newcommand{\TV}{{\mathrm{TV}}} % total variation
\newcommand{\Proj}{{\mathrm{Proj}}}
\newcommand{\prj}{{\mathrm{prj}}}
\newcommand{\prox}{\mathbf{prox}}
\newcommand{\refl}{\mathbf{refl}}
\newcommand{\reflh}{\refl^{\bH}}
\newcommand{\proxh}{\prox^{\bH}}
\newcommand{\minimize}{\text{minimize}}
\newcommand{\bgamma}{\boldsymbol{\gamma}}
\newcommand{\bsigma}{\boldsymbol{\sigma}}
\newcommand{\bomega}{\boldsymbol{\omega}}
\newcommand{\blambda}{\boldsymbol{\lambda}}
\newcommand{\bH}{\vH}
\newcommand{\bbH}{\mathbb{H}}
\newcommand{\bB}{\boldsymbol{\cB}}
\newcommand{\Tau}{\mathrm{T}}
\newcommand{\tnabla}{\widetilde{\nabla}}
\newcommand{\TDRS}{T_{\mathrm{DRS}}}
\newcommand{\TPRS}{T_{\mathrm{PRS}}}
\newcommand{\TFBS}{T_{\mathrm{FBS}}}
\newcommand{\best}{\mathrm{best}}
\newcommand{\kbest}{k_{\best}}
\newcommand{\diff}{\mathrm{diff}}
\newcommand{\barx}{\bar{x}}
\newcommand{\xgbar}{\bar{x}_g}
\newcommand{\xfbar}{\bar{x}_f}
\newcommand{\hatxi}{\hat{\xi}}
\newcommand{\xg}{x_g}
\newcommand{\xf}{x_f}
\newcommand{\du}{\mathrm{d}u}
\newcommand{\dy}{\mathrm{d}y}
\newcommand{\kconvergence}{\stackrel{k \rightarrow \infty}{\rightarrow }}
\DeclareMathOperator{\shrink}{shrink} % shrinkage
\DeclareMathOperator*{\argmin}{arg\,min}
\DeclareMathOperator*{\argmax}{arg\,max}
\DeclareMathOperator*{\Min}{minimize}
\DeclareMathOperator*{\Max}{maximize}
\DeclareMathOperator*{\Fix}{Fix}
\DeclareMathOperator*{\zer}{zer}
\DeclareMathOperator*{\nablah}{\nabla^{\bH}}
\DeclareMathOperator*{\gra}{gra}
\DeclarePairedDelimiter{\dotpb}{\langle}{\rangle_{\bH}}
\DeclarePairedDelimiter{\dotpv}{\langle}{\rangle_{\vH}}
\DeclarePairedDelimiter{\dotp}{\langle}{\rangle}

%% macros for environments math equations

\newcommand{\MyFigure}[1]{../fig/#1}

\newcommand{\bc}{\begin{center}}
\newcommand{\ec}{\end{center}}

\newcommand{\bdm}{\begin{displaymath}}
\newcommand{\edm}{\end{displaymath}}

\newcommand{\beq}{\begin{equation}}
\newcommand{\eeq}{\end{equation}}

\newcommand{\bfl}{\begin{flushleft}}
\newcommand{\efl}{\end{flushleft}}

\newcommand{\bt}{\begin{tabbing}}
\newcommand{\et}{\end{tabbing}}

\newcommand{\beqn}{\begin{align}}
\newcommand{\eeqn}{\end{align}}

\newcommand{\beqs}{\begin{align*}} % no equation numbers
\newcommand{\eeqs}{\end{align*}}  % no equation numbers

%% macros for theorem-like environments

%\newtheorem{theorem}{Theorem}
\newtheorem{assumption}{Assumption}
%\newtheorem{condition}{Condition}
%\newtheorem{rul}{Rule}
%\newtheorem{definition}{Definition}
%\newtheorem{corollary}{Corollary}
%\newtheorem{remark}{Remark}
%\newtheorem{lemma}{Lemma}
%\newtheorem{proposition}{Proposition}
%\newtheorem{example}{Example}
%\newtheorem{proof}{Proof}

%%%%%%%%%%%%%%%%%%%%%%%%%%%%%%%%%%%%%%%%%%%%%

\title{Assignment 2 of \bf{MATP4820}}
\author{Izaak King}
\date{Feb-09-2026}



\begin{document}

\maketitle

\section*{Problem 1}
Consider the sequence $\{x^{(k)}\}_{k=0}^\infty$ defined by
$$x^{(k)}=\left\{\begin{array}{ll}
(\frac{3}{4})^{2^k}, & \text{ if }k\text{ is even}\\[0.1cm]
\frac{2}{k+2}x^{(k-1)}, & \text{ if }k\text{ is odd}
\end{array}\right.$$
\begin{enumerate}
\item
Is this sequence Q-quadratically convergent? If yes, prove it; if not, explain why. 

\vspace{8mm}

No, this sequence is not Q-quadratically convergent. 

First, as $k \rightarrow \infty,$ \[\left(\frac{3}{4}\right)^{2^k} \rightarrow 0, \quad\frac{2}{k+2}x^{(k-1)} \rightarrow 0 \implies x^* = 0\]

Now, if k is even:
\[ \frac{\left\|x^{(k + 1)} - x^* \right\|_2}{\left\|x^{(k)} - x^* \right\|^2_2} = \frac{\left| \frac{2}{(k + 1) + 2}x^{(k + 1 - 1)}\right|}{\left|(\frac{3}{4})^{2^k}\right|^2} = \frac{\left| \frac{2}{k + 3}x^{(k)}\right|}{\left| (\frac{3}{4})^{2^k}\right|^2} = \frac{\left|\frac{2}{k + 3}(\frac{3}{4})^{2^k}\right|}{\left|(\frac{3}{4})^{2^k}\right|^2} = \frac{\left|\frac{2}{k + 3}\right|}{\left|(\frac{3}{4})^{2^k}\right|} = \frac{2}{k + 3}\left(\frac{4}{3}\right)^{2^k}\]

Therefore, as $k \rightarrow \infty$ \[ \frac{\left\|x^{(k + 1)} - x^* \right\|_2}{\left\|x^{(k)} - x^* \right\|^2_2} = \frac{2}{k + 3}\left(\frac{4}{3}\right)^{2^k} \rightarrow \infty\]

Thus, for even k, \[\frac{\left\|x^{(k + 1)} - x^* \right\|_2}{\left\|x^{(k)} - x^* \right\|^2_2} \rightarrow \infty\]

So, there is no M such that \[\frac{\left\|x^{(k + 1)} - x^* \right\|_2}{\left\|x^{(k)} - x^* \right\|^2_2} \leq M,\]

Therefore, $\{x^{(k)}\}_{k=0}^\infty$ is not Q-quadratically convergent

\newpage

\item Is this sequence R-quadratically convergent? If yes, prove it; if not, explain why.
\end{enumerate}

\vspace{8mm}

\noindent Yes the sequence is R-quadratically convergent.

\vspace{4mm}


\begin{proof}First, let $\{x^{(k)}\}_{k=0}^\infty$ converge to $x^*$. To show the sequence $x^{(k)}$ is R-quadratically \begin{adjustwidth}{2em}{}convergent, we have to show there is $\{v^{(k)}\}_{k=0}^\infty$ such that

\[ \left\|x^{(k)} - x^*\right\|_2 \leq v^{(k)}, \text{ and }v^{(k)} \text{ Q-quadratically converges to 0.}\]
Then, as $k \rightarrow \infty$,
\[\left(\frac{3}{4}\right)^{2^k} \rightarrow 0, \text{ and } \frac{2}{k + 2}x^{(k - 1)} = \frac{2}{k + 2}\left(\frac{3}{4}\right)^{2^{(k-1)}} \rightarrow 0 \]
Thus, $x^* = 0$.
Now, take $\{v^{(k)}\}_{k=0}^\infty = \left(\sqrt{\frac{3}{4}}\right)^{2^k}$.

\noindent Then if k is even,
\[ \left\|x^{(k)} - x^*\right\|_2 = \left|\left(\frac{3}{4}\right)^{2^k}\right| \leq \left(\sqrt{\frac{3}{4}}\right)^{2^k} = v^{(k)}\]
If k is odd,
\[ \left\|x^{(k)} - x^*\right\|_2 = \left|\frac{2}{k + 2}x^{(k - 1)}\right| = \left|\frac{2}{k + 2}\left(\frac{3}{4}\right)^{2^{(k - 1)}}\right| \leq \left(\sqrt{\frac{3}{4}}\right)^{2^k} = v^{(k)}\]
Thus, $\left\|x^{(k)} - x^*\right\|_2 \leq v^{(k)}$.

\noindent Now, we will show $\{v^{(k)}\}_{k=0}^\infty$ Q-quadratically converges to 0.

\noindent Note, as $k \rightarrow \infty$, $\left(\sqrt{\frac{3}{4}}\right)^{2^k} \rightarrow 0$, so $v^* = 0$, and $v^{(k)}$ converges to 0.
Then,
\[ \frac{\left\|v^{(k + 1)} - v^*\right\|_2}{\left\|v^{(k)} - v^*\right\|^2_2} = \frac{\left(\sqrt{\frac{3}{4}}\right)^{2^{k + 1}}}{\left(\left(\sqrt{\frac{3}{4}}\right)^{2^k}\right)^2} = 1\]
Since the ratio is bounded by 1, the Q-quadratic condition is satisfied, and $v^{(k)}$ Q-quadratically converges to 0.
\noindent Therefore we have shown there is $\{v^{(k)}\}_{k=0}^\infty$ such that $\left\|x^{(k)} - x^*\right\|_2 \leq v^{(k)}$, and $v^{(k)}$ Q-quadratically converges to 0, so $\{x^{(k)}\}_{k=0}^\infty$ is R-Quadratic.

\end{adjustwidth}

\end{proof}





\newpage

\section*{Problem 2}
Consider the linear system
\begin{equation}\label{eq:lin-sys}
\begin{aligned}
2x_1 - x_2 = 3\\
-x_1 + 4x_2 = 2
\end{aligned}
\end{equation}

Define the iterative update scheme
\begin{equation}\label{eq:scheme-1}
\begin{aligned}
x_1^{(k+1)} = \frac{x_2^{(k)} +3}{2},\\
x_2^{(k+1)} = \frac{x_1^{(k)} + 2}{4}.
\end{aligned}
\end{equation}
Does the sequence $\{(x_1^{(k)},x_2^{(k)})\}_{k=0}^\infty$ generated by the scheme in \eqref{eq:scheme-1} from any initial point $(x_1^{(0)},x_2^{(0)})$ always converge to the solution of the linear system in \eqref{eq:lin-sys}? If yes, prove it; if not, explain why and give a specific initial point $(x_1^{(0)},x_2^{(0)})$ such that the generated sequence converges to the solution of \eqref{eq:lin-sys}.
\vspace{8mm}

\noindent Yes the sequence always converges to the solution of the linear system.

\vspace{4mm}


\begin{proof}First, we solve the linear system,
\begin{adjustwidth}{2em}{}
\begin{equation}\label{eq:lin-sys}
\begin{aligned}
2x_1 - x_2 = 3 \rightarrow x_1 = 2\\
-x_1 + 4x_2 = 2 \rightarrow x_2 = 1
\end{aligned}
\end{equation} 
So, $x^* = 
\begin{bmatrix}
2 \\ 1
\end{bmatrix}
$

\noindent Then, we can take the error of our current sequence $\{x^{(k)}\}_{k = 0}^{\infty} = \{ x_1^{(k)}, x_2^{(k)}\}_{k = 0}^{\infty}$, to be written as $\left\|x^{(k)} - x^*\right\|$. 

\noindent Now, \[ \left\|x^{(k)} - x^*\right\|^{2}_{2} = (x_1^{(k)} - 2)^2 + (x_2^{(k)} - 1)^2\]

\noindent Then, after an iteration, the error of our current solution is given by $\left\|x^{(k + 1)} - x^*\right\|$

\[ \left\|x^{(k + 1)} - x^*\right\|^2_2 = \left\| 
\begin{bmatrix}
\frac{x_2^{(k)} + 3}{2} \\
\frac{x_1^{(k)} + 2}{4} 
\end{bmatrix} - \begin{bmatrix}
2 \\ 
1
\end{bmatrix}
\right\|^2_2 = \left\|
\begin{bmatrix}
\frac{x_2^{(k)}}{2} - \frac{1}{2}\\
\frac{x_1^{(k)}}{4}  - \frac{1}{2}
\end{bmatrix}\right\|^2_2 = \frac{(x_2^{(k)} - 1)^2}{4} + \frac{(x_1^{(k)} - 2)^2}{16}
\]

\noindent Then,
\[ \frac{(x_2^{(k)} - 1)^2}{4} + \frac{(x_1^{(k)} - 2)^2}{16}
\leq \frac{1}{4}\left((x_1^{(k)} - 2)^2 + (x_2^{(k)} - 1)^2\right)\]
\[\implies \left\|x^{(k + 1)} - x^*\right\|^2_2 \leq \frac{1}{4}\left\|x^{(k)} - x^*\right\|^2_2 \implies \left\|x^{(k + 1)} - x^*\right\|_2 \leq \frac{1}{2}\left\|x^{(k)} - x^*\right\|_2\]

\noindent Now, in general, when iterating
\[\left\|x^{(k)} - x^*\right\|_2 \leq \frac{1}{2}^k\left\|x^{(0)} - x^*\right\|_2 \]

\noindent And, as $k \rightarrow \infty$,
\[ \frac{1}{2}^k\left\|x^{(0)} - x^*\right\|_2 \rightarrow 0.\]

\noindent Thus, as we iterate, the error of our solution goes to 0 for any initial point, proving it always converges to the exact solution of the linear system.
\end{adjustwidth}

\end{proof}




\newpage

\section*{Problem 3}
Let $\{x^{(k)}\}_{k=1}^\infty$ and $\{y^{(k)}\}_{k=1}^\infty$ be two scalar sequences. Suppose $x^{(k)}$ Q-sublinearly converges to $x^*$ and $y^{(k)}$ Q-linearly converges to $y^*$. 
\begin{enumerate}
\item Must $x^{(k)} + y^{(k)}$ be Q-sublinearly convergent to $x^*+y^*$? If yes, prove it; if not, give a counter example. 

\vspace{8mm}

\noindent Yes, $x^{(k)} + y^{(k)}$ must be Q-sublinearly convergent to $x^* + y^*$.

\vspace{4mm}


\begin{proof}Assume $\{x^{(k)}\}_{k=1}^\infty$ and $\{y^{(k)}\}_{k=1}^\infty$ are two scalar sequences, and $\{x^{(k)}\}_{k=1}^\infty$
\begin{adjustwidth}{2em}{}
\noindent Q-sublinearly converges to $x^*$, and $\{y^{(k)}\}_{k=1}^\infty$ Q-linearly converges $y^*$

\noindent That is, $k \rightarrow \infty$
\[ \frac{\left\|x^{(k + 1)} - x^*\right\|_2}{\left\|x^{(k)} - x^*\right\|_2} \rightarrow 1 \text{ and } \frac{\left|y^{(k + 1)} - y^*\right\|_2}{\left|y^{(k)} - y^*\right\|_2} \leq r, \text{ for } r\in(0,1)\]

\noindent Now, define $z^{k} = x^{(k)} + y^{(k)}$, $z^* = x^* + y^*$. We will show $z^{(k)}$ is Q-sublinearly convergent.

\noindent First note, as $k \rightarrow \infty$
\[ \frac{\left\|x^{(k + 1)} - x^*\right\|_2}{\left\|x^{(k)} - x^*\right\|_2} \rightarrow 1 \text{ and } \frac{\left\|y^{(k + 1)} - y^*\right\|_2}{\left\|y^{(k)} - y^*\right\|_2} \leq r, \text{ for } r\in(0,1)\]

\noindent This implies $\frac{\left\|y^{(k)} - y^*\right\|_2}{\left\|x^{(k)} - x^*\right\|_2}\rightarrow 0$, because the Q-sublinear sequence $x^{(k)}$ decays slower than the Q-linear sequence $y^{(k)}$.


\noindent Now, 

\[ \frac{\left\|z^{(k + 1)} - z^*\right\|_2}{\left\|z^{(k)} - z^*\right\|_2} = \frac{\left|x^{(k + 1)} - x^* + y^{(k + 1)} - y^*\right|}{\left|x^{(k)} - x^* + y^{(k)} - y^*\right|} = \frac{\left|\frac{x^{(k + 1)} - x^*}{x^{(k)} - x^*} + \frac{y^{(k + 1)} - y^*}{x^{(k)} - x^*}\right|}{\left| 1 + \frac{y^{(k)} - y^*}{x^{(k)} - x^*}\right|}\]

\noindent Now, as we have shown, as $k \rightarrow \infty$,

\[ \frac{x^{(k + 1)} - x^*}{x^{(k)} - x^*} \rightarrow 1, \quad \frac{y^{(k + 1)} - y^*}{x^{(k)} - x^*} \rightarrow 0, \quad \frac{y^{(k)} - y^*}{x^{(k)} - x^*} \rightarrow 0,\]

\noindent Thus, as $k \rightarrow \infty$,
\[ \frac{\left\|z^{(k + 1)} - z^*\right\|_2}{\left\|z^{(k)} - z^*\right\|_2} = \frac{\left|\frac{x^{(k + 1)} - x^*}{x^{(k)} - x^*} + \frac{y^{(k + 1)} - y^*}{x^{(k)} - x^*}\right|}{\left| 1 + \frac{y^{(k)} - y^*}{x^{(k)} - x^*}\right|} \rightarrow 1\]

\noindent So, $z^{(k)}$ is Q-sublinearly convergent. Since $z^{(k)}$ was defined as $z^{(k)} = x^{(k)} + y^{(k)}$ with $x^{(k)}$ Q-sublinear convergent and $y^{(k)}$ Q-linear convergent, we have proven our claim.


\end{adjustwidth}

\end{proof}


\newpage


\item Can $x^{(k)} + y^{(k)}$ be Q-linearly convergent to $x^*+y^*$? If yes, give an example; if not, explain why.
\vspace{8mm}

No, $x^{(k)} + y^{(k)}$ can never be Q-linearly convergent.

\vspace{4mm}


As proven in Part 1, define $z^{(k)} = x^{(k)} + y^{(k)}$ and $z^* = x^* + y^*$. Then
\[
\frac{\left\|z^{(k+1)} - z^*\right\|_2}{\left\|z^{(k)} - z^*\right\|_2} \rightarrow 1 \quad \text{as } k \rightarrow \infty.
\]

By definition, a sequence is Q-linearly convergent if there exists $r \in (0,1)$ such that
\[
\frac{\left\|z^{(k+1)} - z^*\right\|_2}{\left\|z^{(k)} - z^*\right\|_2} \leq r \quad \text{for all sufficiently large } k.
\]

Since the limit of the ratio for $z^{(k)}$ is 1, which is not less than 1, it is impossible for the sequence to satisfy the Q-linear convergence criterion. Therefore, $x^{(k)} + y^{(k)}$ cannot be Q-linearly convergent. 

\vspace{2mm}

\end{enumerate}

\end{document}
